\documentclass{beamer}
\usetheme{Madrid}
\usepackage[utf8]{inputenc}
\usepackage[T2A]{fontenc}
\usepackage[russian]{babel}
\usepackage{graphicx, booktabs, array}
\setbeamercolor{background canvas}{bg=blue!25}
\usebeamercolor*{normal text}

\title{История Орловского государственного\\ университета имени И. С. Тургенева}
\author{Свиридов Д.Д.}
\institute{Орловский государственный университет имени И. С. Тургенева}
\date{25 мая 2026}

\setbeamertemplate{footline}{
  \leavevmode
  \hbox{
  \begin{beamercolorbox}[wd=.333333\paperwidth,ht=2.25ex,dp=1ex,center]{author in head/foot}%
    \usebeamerfont{author in head/foot}\insertshortauthor
  \end{beamercolorbox}
  \begin{beamercolorbox}[wd=.333333\paperwidth,ht=2.25ex,dp=1ex,center]{title in head/foot}%
    \usebeamerfont{title in head/foot}\insertshorttitle
  \end{beamercolorbox}
  \begin{beamercolorbox}[wd=.333333\paperwidth,ht=2.25ex,dp=1ex,right]{date in head/foot}%
    \usebeamerfont{date in head/foot}\insertshortdate{}\hspace*{2em}
    \insertframenumber{} / \inserttotalframenumber\hspace*{2ex} 
  \end{beamercolorbox}}
  \vskip0pt
}

\setbeamertemplate{navigation symbols}{}
\setbeamercovered{transparent}

\begin{document}

\begin{frame}
  \titlepage
\end{frame}

\begin{frame}{Основание университета}
  \begin{columns}[T]
    \begin{column}{0.6\textwidth}
      \begin{itemize}
        \item Дата основания - \textbf{1919 год}.
        \item Первое название - «Орловский педагогический институт».
        \item 1931 год - новое название: индустриально-педагогический институт.
        \item Университетский статус - с 1996 года.
        \item Имя И. С. Тургенева - с 2010 года.
      \end{itemize}
    \end{column}
    \begin{column}{0.4\textwidth}
      \includegraphics[width=\textwidth]{ахаха.jpg}
    \end{column}
  \end{columns}
\end{frame}

\begin{frame}{Развитие в советский период}
  \begin{itemize}
    \item 1920-1930-е годы - статус крупного педагогического центра.
    \item 1941-1943 годы - разрушение зданий и пауза в учебном процессе.
    \item 1943 год - начало послевоенного восстановления.
    \item Конец 1960-х годов - 7 факультетов и более 3000 студентов.
  \end{itemize}
  \vfill
  \begin{block}{Важный факт}
    1972 год - орден «Знак Почёта» за вклад в подготовку кадров.
  \end{block}
\end{frame}

\begin{frame}{Ключевые вехи истории}
  \begin{table}
    \centering
    \begin{tabular}{p{0.25\textwidth} p{0.65\textwidth}}
      \toprule
      \textbf{Год} & \textbf{Событие} \\
      \midrule
      1919 & Основание Орловского педагогического института \\
      1931 & Преобразование в индустриально-педагогический институт \\
      1941–1943 & Эвакуация и разрушения в годы войны \\
      1972 & Награждение орденом «Знак Почёта» \\
      1996 & Получение статуса классического университета \\
      2010 & Присвоение имени И. С. Тургенева \\
      \bottomrule
    \end{tabular}
  \end{table}
\end{frame}

\begin{frame}{Современный университет (после 2000 года)}
  \begin{columns}
    \begin{column}{0.5\textwidth}
      \begin{itemize}
          \item 12 институтов и факультетов;
          \item более 10 000 студентов;
          \item 50 и более направлений подготовки.
      \end{itemize}
    \end{column}
    \begin{column}{0.45\textwidth}
      \includegraphics[width=\textwidth]{огу ы.jpg}
    \end{column}
  \end{columns}
\end{frame}


\begin{frame}{Роль в науке и образовании региона}
  \begin{block}{Научные достижения}
    \begin{itemize}
      \item Первое место в Орловской области по объёму НИОКР.
      \item Диссертационные советы: педагогика, физика, экономика, филология.
      \item Ежегодная международная конференция - «Тургеневские чтения».
    \end{itemize}
  \end{block}
  \vfill
  \begin{exampleblock}{Социальная и культурная миссия}
    Музей Тургенева, студенческие театры, волонтёрские проекты, культурные мероприятия.
  \end{exampleblock}
  \vfill
  \centering
  \textit{ОГУ им. И.С. Тургенева - гордость Орла.}
\end{frame}

\end{document}